% Generated draft: requires mathematical and source review.
\documentclass[11pt]{article}
\usepackage[margin=25mm]{geometry}
\usepackage{fontspec}
\setmainfont{Latin Modern Roman}
\usepackage{amsmath,amssymb,mathtools,unicode-math}
\setmathfont{Latin Modern Math}
\usepackage{enumitem,needspace,xurl,hyperref}
\setlength{\parindent}{0pt}
\setlength{\parskip}{5pt}
\setlength{\emergencystretch}{3em}
\newcommand{\N}{\mathbb{N}}
\newcommand{\Z}{\mathbb{Z}}
\newcommand{\Q}{\mathbb{Q}}
\newcommand{\R}{\mathbb{R}}
\newcommand{\C}{\mathbb{C}}
\newcommand{\F}{\mathbb{F}}
\newcommand{\A}{\mathbb{A}}
\newcommand{\PP}{\mathbb{P}}
\newcommand{\E}{\mathbb{E}}
\newcommand{\eps}{\varepsilon}
\newcommand{\dd}{\,\mathrm{d}}
\begin{document}
\section*{UnsolvedMath ID: 30000617}
\textbf{Source problem code:} catalogue entry
\subsection{Definitions and mathematical statement}
Let $\varepsilon_0$ be the least ordinal $\alpha>0$ satisfying $\omega^\alpha=\alpha$. Let $N$ denote the intended additive norm on codes for ordinals below $\varepsilon_0$, and let $F$ be the intended numerical control function. The source asks for a transition between provability and unprovability of its controlled finite-well-ordering assertions $\mathrm{FWO}_N(F)$ as the growth of $F$ varies.

Reading ``descending pair'' literally as strict ordinal descent, the indicated quantifier pattern would be
\[
\begin{split}
\forall K\in\N\ \exists m\geq1\ \forall(\alpha_0,\ldots,\alpha_{m-1})\in\varepsilon_0^m,
\quad&
\left[\forall i<m,\ N(\alpha_i)\leq K+F(i)\right]\\
&\Longrightarrow\left[\exists i<j<m,\ \alpha_i>\alpha_j\right].
\end{split}
\]
This literal formulation cannot express a usual finite-well-ordering principle: whenever the bound admits $0$, the constant zero sequence has no strictly descending pair, regardless of length. The problem is therefore not repaired by merely choosing a longer sequence.

A usual well-ordering formulation would instead require a nondecreasing pair $\alpha_i\leq\alpha_j$, or bound the lengths of controlled strictly decreasing sequences. These are possible intended corrections, not verified replacements of the source's ``descending pair''. The extract also gives no recursive definition of the additive norm, no ordinal coding, no parameterized family of functions $F$, and no formal arithmetic theory in which provability is measured. Those data and the order direction must be recovered before a precise phase-transition conjecture can be stated. No critical growth rate or independence theorem is inferred.
\subsection{Short English statement}
Determine the proof-theoretic growth threshold for controlled finite well-ordering below $\varepsilon_0$. The extract's literal descending-pair requirement fails for constant sequences, and the norm, function family and base theory are missing, so the intended phase-transition assertion remains unresolved.
\subsection{Sources}
\begin{itemize}
\item Supplied report: \emph{Mini-workshop: Logic, Combinatorics and Independence Results} (2007), Oberwolfach Reports, DOI 10.4171/owr/2006/52, pp. 3087--3130; proposed year 2006. \url{https://doi.org/10.4171/owr/2006/52}.
\item Local UnsolvedMath registry, permanent ID 30000617, catalogue entry; created 25 August 2026, updated 16 September 2026. It labels the statement corrected\_verified after comparison with the official report. Its supplied literature triage of 21 August 2026 labels the extracted claim open and reports no located authoritative complete resolution. These source labels do not resolve the direction and definition gaps identified above. No further external reference is supplied, and no independent literature assessment is asserted.
\end{itemize}
\end{document}
